---
format: yohaku/0.2
target: text
macro: '[[macros/raw]]'
---

## 正文
````text
# π 研究中的工作流形成与跨任务记忆：v03–v06c 观察

记录日期：2026-11-12。对象：Yohaku 中同一名 `adaption / glm-53-flash / none` π 研究负责人。本文根据封存过程包及独立核对记录撰写，现已覆盖 v05 的真实运行和 v06/v06c 的模型请求前启动失败。v05 在 12/16 次请求后由负责人主动结束，同根累计 38/42；全部旧 79 笔未知保留、新增未知为 0，阶段已关闭。此前的截断、失败、离线修复与首次收尾超时均保留。

目前的证据已从“方法送达和写入记忆”推进到“跨上下文选择旧脚本、实际执行、读取结果并保存独立说明文档”。v05 的新目标没有指定操作步骤，但旧记忆明确保留了验证待办，外部审查也影响了选择，因此这是受历史引导的自主接续与复用案例，尚不是新工作流的独立发明。截至 v05 未创建辅助者；v06c 由宿主创建了一个尚未执行任务的 Sol 节点，仍未发生实际联合研究。稳定分化、参数训练或方法的因果收益尚未建立。

## 实验问题与设置

研究载体是既有 π 逼近项目。目标包括观察 AI 如何读取计算结果、留下方法、处理纠错，以及在后续任务中是否采用此前形成的方法。研究者选择具体工具和处理步骤；宿主提供工具、文件、反馈、真实请求额度和导出机制。

v04 的 A、B 是同一负责人的两个任务上下文，不是两名独立辅助 AI。A 接受明确的外部数学审查，决定如何核验和修订；B 以新上下文处理一个可验证的后续目标：从既有精确有理数给出小数点后 18 位的有保证十进制上界，并说明其数学含义。B 初始输入没有 A 的对话、预期答案或精确方法路径，正常持久记忆与项目文件可按需读取。

这是受到外部条件约束的案例观察。显式反馈、要求保存方法、后续任务和额度均由外部设置；不能把全部过程描述为无提示、自发演化。

## 阶段与实际请求

这里的 POST 是模型请求次数，不是研究时长，也不是计算工具调用次数。

| 阶段 | 实际发送模型 POST | 完整响应 | 新增未知 | 主要可见行动 |
|---|---:|---:|---:|---|
| v03 | 8 | 7 | 1 | 修改计算文件，完成一次计算；流式请求中断后保留未知预留 |
| v03b | 4 | 4 | 0 | 保存方法，修订记忆与研究进展；原阶段累计到 12/12 |
| v04 A | 5 | 5 | 0 | 读取外部审查，保存回应并修订三份研究文件 |
| v04 B | 3 | 3 | 0 | 读取修订记忆、结果索引与目录，主动请求方法和计算源 |
| v04 离线修复 | 0 | 不适用 | 0 | 隔离原生创建和登记闭环、回执回放、封存 B 检查点 |
| v04c B 接续 | 6 | 6 | 0 | 收到方法并用于结果读取和保存记忆；项目清单错误阻止复核运行与独立文档保存 |
| v05 自主研究 | 12 | 12 | 0 | 自行接续既有验证待办，执行旧脚本、读取报告、保存说明并修订正文；主动结束，未用满 16 次上限 |
| v06 首次启动 | 0 | 不适用 | 0 | 创建混合模型同伴时被旧 profile 校验拒绝，尚未请求模型 |
| v06c 人工接续 | 0 | 不适用 | 0 | 宿主创建一个空闲 Sol 节点；负责人入口因等待工具契约不符终止，Sol 原生 turn 亦为 0 |

v04 及之前累计 20 次实际 POST 的旧阶段保持关闭；v04c 新增 6 次，累计 26/26，现也已关闭。历史未知记录合计 79 笔且逐值保留。截至 v04c，新增辅助者与求助申请均为 0。v04c 的辅助者上限明确为 0，因此本轮不能用于判断模型是否倾向自主分工或创建同伴。此前各阶段的零辅助者记录也不单独构成协作机制失败。

v05 另经用户授权共享最多 16 次新 POST、3.0408704 Credits 最坏责任和 2 个可选辅助者（深度 1，无扩张权限）。实际使用 12 次，责任保留 2.2806528 Credits；这是责任上界，不是平台实扣。负责人最终无新工具调用而结束，触发共享停止。剩余 4 次随阶段关闭而未使用，不能视作下一轮自动授权。v05 仍为 0 次求助、0 个辅助者；它没有发生实际分工。

## 观察一：修订记忆确实进入了新任务

A 面对外部指出的 `inf` / `liminf` 混淆，读取现有定义和项目记录，接受核心异议，保存回应，修订进展、方法和记忆。随后 B 的实际已发送请求中出现了这份新记忆。该结论来自请求输入与文件版本的核对，而不只是“文件已保存”或界面提示。

修订记忆版本为 `4a43a7331fc6772ca0773edf4a529dfee58a8938ffffd22c58fe4e048225adba`。此证据支持跨任务信息持久化与交付；不能推出稳定的专业分化，更不能推出模型参数发生训练。

数学纠错来自明确的外部反馈。核心推断得到修订，但回应仍有局部公式与下标问题，另存为待核验的外部意见；回应保存或登记均不构成数学批准。

## 观察二：v04 的主动检索与未交付终止点

原 v04 B 的交付过程如下；这里保留其当时被截断的历史：

| 模型轮次 | 实际收到的内容 | 随后选择的行动 |
|---|---|---|
| 0 | 新目标与通常项目入口 | 读取记忆、列出项目目录 |
| 1 | 修订记忆与项目目录 | 读取结果索引、列出研究目录 |
| 2 | 结果预览与研究目录 | 请求方法文件和计算源 |
| 3 | 请求在网络前因额度耗尽被拒绝 | 方法与计算源未送入模型 |

方法文件版本为 `64c95fcec6039eedb7b7e80607a8f24c0c8574cf0b80d336005e8868ee4d248b`。两次读取已成功返回宿主，但只装入了未发送的下一轮输入。B 没有在本轮保存目标产物，也没有可见的计算调用。

因此必须分别记录三个事件：工具返回宿主、内容进入真实模型请求、后续行动实际采用内容。这三个事件不可互换。仅有方法名称引用或最终答案正确，也不足以单独证明方法产生了因果效果。

## 观察三：额度既约束资源，也截断了观察

真实请求入口拒绝了超额请求，换任务或子授权不会自动获得新额度。未知请求保留责任预留，没有通过重放 POST 消除不确定状态。

但 B 的三次生成主要用于寻找资料，最后选择的工具返回尚未交付，观察就结束了。这是预算造成的截断，不能据此推断模型没有复用能力。后续实验应预留接收工具结果、处理、保存以及终态所需的请求，同时继续保留总量硬上限。

## 观察四：产物保存、登记与批准需要独立证据

A 的四份文件确实保存成功，但三次正式作者回应登记全部失败。原因是原生创建回执缺少观察器预期的版本字段。后续一版通过 18 项合成检查的修复，仍遗漏了真实源文件尾部的完整输出路由段，独立回放原始回执才发现不兼容。

最终隔离修复比较完整原生格式、正文、路由、任务、版本快照与生成物，另外公开标注宿主的 `host_write_evidence`，没有伪造原生 `revision`。临时原生闭环产生 1 条作者回应和 1 条明确版本继承，状态仍为 `pending_independent_review`，独立批准为 0。完整闭环由导出脚本、原生进程日志、文件与数据库相互佐证；独立复查另行回放了归档回执及篡改检查。

截至 v04 离线修复时，修复只在隔离候选验证，原研究运行候选未切换。三次历史登记失败被保留，正文中提前写的“已登记”由宿主状态冲突提示标出。不能用临时成功补写历史成功，也不能让工程状态替代数学判断。

## 观察五：v04c 收到方法后出现了实际采用，也暴露了配置自锁

v04c 是关联原 B 的受控接续。首个真实请求保留原来 18 项输入的逐值相同前缀，两笔待交付工具结果仍各出现一次、顺序不变；仅追加额度与检查点接续的操作通知。两项正文实际进入已完成的模型轮次。这与原 v04 的“工具已返回、模型未收到”是两个不同终止点。

接续中的可见行为包括：按结构化结果接口读取原产物的 `runA/min_markov_upper_bound/decimal_40`；结合计算源中的向下截断约定处理 18 位显示；将来源哈希、下截断与向上一单位的规则、有限见证的适用范围写入持久记忆。这些实际读取与保存内容支持“本次方法被采用”，证据超出了引用方法名称。

保存记忆的源版本为 `7ba7b6d1bfa8ef67bb8cd99663fc8a9e6e506a597eb81773d9e1097fb199d2d6`。其中的 18 位十进制上界 `0.000048130561246290`，经独立整数运算检查，确为原存储有理数的严格上界及最紧 18 位向上表示。检查范围是该显示与原存储分数的关系，并不把整份记忆、π 证书或全部数学论证标为通过。记忆仍带有旧登记状态、旧下标及其他未核验或错误陈述，必须按具体主张分别审查。

负责人还保存了复核脚本，但没有成功提交或执行；独立结果文档也未保存成功。主要阻塞来自项目清单：工具接受并保存了含大写字母的动作键 `verify-runA-dec18`，而项目清单要求动作名为小写 ASCII、数字、`-` 或 `_`。随后项目作用域的计划、物化、读取、写入乃至修复操作均在清单验证处返回 `invalid_manifest`。保存入口允许非法配置落地，而后续工具又依赖整份配置有效，形成了自锁。

模型的修复尝试没有解除该状态。六次完整响应及下一次预算拒绝之后，宿主仅将动作键改为 `verify-runa-dec18`，恢复清单可验证性；脚本、参数、输出路径和数学正文未因这次宿主操作改变，也未运行研究代码。这是明确的外部工程干预，不能记为模型自恢复或数学验证成功。

这一轮完成了方法交付和一次可见采用，但交付物不完整、工程阻塞需要宿主介入。它提供了更具体的工作流证据，也保留了失败的边界。

## 观察六：v05 离线准备暴露了停止契约与实现的不一致

v05 准备包记录新增模型 POST 为 0，原 26/26 阶段关闭，79 笔旧未知保留，正式候选未切换。拟议下一阶段为所有参与者共享 16 次新 POST、最多 3.0408704 Credits 最坏责任、2 个可选辅助者且深度为 1；同根总上限为 42。这些仍是待授权方案，不是已经发生的运行。

独立审阅核对了 194 份清单文件、20 份改动源码及包装的指纹。归档报告配置事务 31 项、接续 13 项、额度 11 项检查通过；本次复查没有把这些数目记为重新执行通过。配置部分的 128 次原生调用、44 次公开工具调用和 3 个有效快照相互吻合；当前研究节点可达的候选预验证、CAS、有效配置切换与受限修复路径未发现启动阻塞。四项旧调用及待交付结果与 v04c 逐值一致、原序各一次；新目标和普通记忆吻合，观察契约未进入输入预览。冻结包不能替代启动时对实际机器状态的复核。

但停止条件存在明确反例。方案要求新增未知出现后停止整个阶段的新请求，而候选 `Gate.note_outcome(..., "unknown", ...)` 只更新当前任务，阶段仍然开放；`post_locked` 没有检查共享停止信号。归档测试 `test_unknown_retained_and_seventeenth_blocked` 甚至在首次 SSL 未知后继续发送 15 次，把总量不超 16 当作通过条件。这验证了总量上限，却没有实现所承诺的提前停止条件。

审阅者使用实际候选 Gate、临时合成授权和历史状态做了独立顺序复现：保留 79 笔旧未知；先确认正常在途期间其他参与者可以准入；再明确登记负责人本次为 unknown；已启动辅助者的下一次 POST 准入依然通过，阶段仍为 open。复现没有网络调用或真实模型请求，也没有修改生产文件；其范围是确定性准入反例，不是完整并发运行验收。

启动前需要持久化的全阶段停止信号，并与新 POST 准入在同一权威锁下判断。必须区分旧未知、正常在途请求和新确认未知：底层 `reserved_outcome_unknown` 在发送前预留时就会写入，不能据此直接把所有在途请求当作终态未知。未知的次数和费用责任继续保留；已在途请求仍需记录回执或清理，不能把本地停止写成远端取消。

另外，负责人 done 也只更新单个任务；宿主执行器的停止新派发、参与者清理与阶段关闭需在激活前补齐执行证据。同一负责人的新上下文段必须清除旧会话链并保留历史，13 项夹具直接使用空 session，尚不能替代真实派发绑定验证。另一个已独立复现的注册表读写校验不一致，可由通用宿主写入触发；当前研究节点无权写该路径，因此列为平台推广前的非本轮阻塞项。

本轮结论是准备的大部分证据成立，当前候选仍应先修复停止缺口，再申请启动。测试通过的数量不能替代对目标行为的核对。本轮没有新的数学计算或自主工作流运行证据。

## 观察七：共享停止反例已修复，监督循环仍需处理正常的申请失效

后续修复包仍为 0 次新模型 POST，原 26/26 和 79 笔旧未知保留。独立核对 134 份清单文件和 12 份变更源码/包装，测试引用的最终源码均可追溯；未改变的预算守卫与任务运行器也与既有归档指纹匹配。归档分别报告共享停止 26 项、正常上下文 18 项、宿主监督 9 项离线验证通过，夹具拥有进程残留为 0。这些是分组的原测试数量，不是本次独立重跑数量。

原“确认新未知后辅助者仍获准”的反例已在新 Gate 上独立复核消除。另行执行的 11 个定向边界检查通过，覆盖旧未知/正常在途不误停、跨参与者停止、排队启动、重启、晚到回执、日志崩溃窗口、负责人结束及共享额度。它们使用实际 Gate 与原预算模块的锁和原子写函数、合成私有状态，未发送网络或模型请求；不是完整并发运行验收。停止标记持久化在共享账本，预留和发送许可分别检查；许可不等于实际已上网。

正常新上下文入口也已补齐：从非空旧 session 出发，先归档原节点、私有任务和记忆，再通过正常入口清除会话链、提升代次并派发。归档集成测试实际调用 Control.begin_context_segment 与 Control.node_dispatch，仅替换付费 worker；其 55 次原生注册表进程均返回 0。四项旧反馈仍保留，输入捕获不被计作实际模型送达。

但新宿主监督程序还有一个可达异常路径。负责人提交带文件版本引用的有效求助后，继续研究并更新了该文件；宿主审批时版本已过期，Gate.approve_pending 抛出 GateError。Supervisor.tick 没有将其转为申请拒绝，而 run 的异常保护只覆盖准备与首次派发，不覆盖后续 tick 循环，因而监督进程可以退出，阶段却仍开放。

独立最小复现使用实际 Gate.request、Supervisor.tick 和后续 POST 准入方法。在尚有辅助名额的合成状态中，过期引用使申请继续处于 pending、共享停止为空，另一名辅助者仍获准提出下一次 POST。另行检查真实 run 方法确认 tick 异常直接向外传播，未触发停止或清理。全程 0 网络请求，未启动真实模型或修改生产文件；此证据的范围是异常路径，不是现场故障报告。

这项修复应保持局部：预期的文件版本失效成为有理由的单项申请拒绝，研究负责人决定是否重提；监督循环遇到无法继续的异常，则先持久关闭全阶段新准入，再进行有界收尾，无法确认清理时保留明确的 cleanup_incomplete。须保留原错误及状态，不把进程退出写成阶段关闭。当前仍不建议激活模型额度；原定 16 次共享 POST、3.0408704 Credits、2 个辅助者、同根 42 次的方案没有扩大。

## 观察八：申请局部拒绝与监督异常收尾通过定向复核

最后的审批与异常处理差量包仍为 0 次新模型 POST，未激活新阶段或额度。独立核对 57 份清单文件、两份修改模块及其源包装、7 份最终报告和 24 个被测源码引用指纹；相对于上一封存候选仅变更 Gate 与 Supervisor。原预算、授权和共享停止的关键方法结构保持不变，既有新上下文及清理实现按原指纹复用。

Gate 新增窄异常类型，将当前源缺失或引用版本失效转为该申请的持久拒绝回执，保留原 request、原引用、预期/实际版本及原因。独立复现确认同批有效申请仍获批准，下一审批回合继续；重复审批不重复拒绝、不改写引用、不替负责人重提。意外访问故障仍向监督器传播，避免将系统故障伪装成普通申请失败。

Supervisor 的定向独立复核使用实际 Gate 与实际 tick/run/故障记录路径、临时合成私有状态和原本地锁。原先过期申请导致监督退出的反例不再出现，连续两个 tick 均保持 monitoring。注入意外审批异常后，先持久化 SESSION 共享停止，再写独立故障记录并进入有界清理；辅助者新 POST 被拒绝。清理无法确认时保留 cleanup_incomplete，重启的 prepare/dispatch 调用数为 0，不重试原坏审批或重放请求。

共享停止后端不可写也被明确区分：两次有界写入失败后记录 admission_stop_confirmed=false，不谎称全阶段已成功停止或完成清理；本监督器的私有故障记录阻止重启重派。后端恢复后继续停止和原收尾。这个结果不构成对不可写存储的全局停止保证，保留相应证据限制。

归档的新增验证分别为审批 10 项、监督联动 1 项、异常 8 项、控制流 2 项，夹具拥有进程残留为 0。审阅者另做上述无网络反例复核，未重新运行上传的宿主脚本、原生夹具或实际模型；合成清理未确认的情况明确保持未完成，没有冒充真实进程清理成功。

本次准备审查在原范围内未发现剩余启动阻塞，结论为可进入已明确授权的有界试跑。最终候选 manifest SHA-256 为 `ea19bf1433ae7ae6a228e3f13607d16914f8ab51703ba1d8ebc1175151503b86`。仍需由用户授权共享最多 16 次新 POST、3.0408704 Credits 最坏责任、最多 2 个深度 1 辅助者、同根累计最多 42 次；启动时核对实际价格和策略绑定，旧 26/26 与 79 笔未知保持。审查结论本身不激活额度，也不证明供应商稳定性或自主研究成果。

## 观察九：开放目标下，旧验证计划完成了执行与留档

v05 真实首请求使用用户的开放目标原文：“继续原来的 π 研究，自行判断下一步值得解决的问题，并组织完成它所需的工作。遇到能力或资源缺口，可以提出具体诉求和方案；把有用的成果与经验留下。”旧 session 经正常入口归档，新上下文绑定原负责人。四项 v04c 未交付工具回执与原调用逐对象一致，在首请求中各一份且保持原序，没有重新执行旧工具。后续请求使用累积输入，旧回执会随上下文再次携带；这里的“各一次”不表示整个网络过程只传输一次。

负责人正常记忆里已有修复后运行验证脚本的待办。它还读取了外部数学与登记状态审查卡。因此可以观察它怎样选择和组织下一步，但不能把沿用旧计划、响应外部意见写成从零发明流程。实际顺序为检查项目与计算接口、修订回应、物化项目并提交已有动作、读取执行结果、保存独立结果文档、更新进展和记忆。全程 32 项新工具调用，唯一新工程任务运行原 `verify_runA_dec18.py`；其字节与 v04c 保存版本相同。

该任务有真实退出码 0、stdout 和报告产物。报告 SHA-256 为 `51d8b8cf1d9bad153f8c9023503e1522f5d1ed3ab9f64bbde82c0270ebb494ff`。负责人随后实际读取结果，并保存 `RUNA-DEC18-BOUND.md.tex`。这补上了 v04c 的“脚本已保存但未执行、说明文档创建失败”两个缺口。脚本其中一路调用原计算实现，另一路仅重写连分数部分而共用 π 区间构造；不能据此称为完全独立的数学认证。

本次外部审阅另写自包含整数检查器，使用 Machin 公式和向外取整的反正切级数，直接认证 runA 的有限见证及其 40 位夹逼。这个检查是审阅者做的，不能归给研究负责人；完整数学边界见配套数学笔记。本轮研究文档虽修正了下标和分母位数，仍把归一化质量漏写一个 `q_n` 因子，并保留错误的等值说明。登记成功与脚本成功不能推出正文每项数学陈述正确。

## 观察十：局部恢复、状态冲突与自述偏差均可从行动核对

新工具调用中有 3 次失败：两次回应登记遇到 `Ancestry already bound to another parent`，一次创建文档时把目录放进文件名。负责人下一轮自行拆开目录与文件名后保存成功，这是本轮可见的局部工具错误恢复；两次登记冲突仍保留，宿主没有绕过。七次登记调用中五次成功，状态均为 `pending_independent_review`，数学字段为 `not_determined_by_host`。

模型最终自述把上轮清单补丁的反馈说成了本轮首轮新动作，并给出与原错误不符的失败原因。实际 32 次新调用没有这项清单补丁；它是首请求补交的历史回执。模型还将责任额度近似写为 2.09，而权威账本为 2.2806528。原模型文字保留，事实差异另行标注。跨上下文携带历史反馈有助于恢复，但事件归属与计量仍须依据工具日志和账本。

首次有界收尾约 123 秒后留下 `cleanup_incomplete`。原 group watcher 约 170 秒后自然完成；宿主从同一共享停止状态进入同一 supervisor 的收尾分支，最终关闭本阶段。此过程没有新上下文、重新派发或新增 POST。最终观察确认研究、工程及收尾进程退出并释放运行租约；阶段关闭与最后一个收尾 watcher 退出不是同一时刻。收尾慢操作的根因仍未定位，不能写成已修复。

## 对分工与分化机制的含义

这些观察给设计带来几个可检验的要求，而不是对生物机制的实验证明：

- 共享规则可以规定读取、写入、预算、求助和权限；具体方法与专长通过任务中的版本化产物和记忆保留。
- 持久化不是最后一步。方法还需要经过检索、实际交付、采用与结果检查，才能建立可复用能力的证据。
- 局部子系统应保存自己的状态与需求，只在需要时向外发出明确请求，减少把全部历史广播给所有模型的成本。
- 扩展需要绑定原任务、具体诉求、交付物、检查方式与额度。辅助者不自动继承再次扩展的权限；原负责人仍承担整合与验收责任。
- 预算拒绝、工具失败和未交付返回应作为实验数据保存。无法区分这些状态时，容易把工程失败误判为模型选择。
- 项目配置应先在候选状态完成校验，再原子提交；保留最后有效版本，并提供不依赖当前损坏清单的修复入口。
- 方法记录应携带适用条件以及成功、失败和外部干预履历，避免把一次部分成功固化成普遍规则。
- 后续应以真实变式任务检查迁移，例如改变精度、对象或证据缺失条件，并分别记录交付、采用和正确性。

当前证据覆盖“修订和跨任务记忆交付”“自主方法检索”、v04c 中的方法采用，以及 v05 中跨上下文继续既有计划、执行检查并保存交付物。尚未观察到新流程发明、跨参与者分工或方法迁移。后续值得观察它在新的真实问题中是否继续使用、修改或放弃旧方法；不应为了得到分工现象而强制创建同伴。长期稳定分化、可推广的效率收益、因果效果与参数训练仍未建立。

数学结果的独立边界说明见 [π 连分数计算审计](https://github.com/epsilon-lain/lain-s-note/tree/main/research/yohaku-pi-cf-audit-2026-11-12)。本记录与仓库中较早的正式多代理试验是不同案例，次数不合并。

## 证据与复核边界

机器可读摘要见 [observations.json](observations.json)，包含来源过程包名称、SHA-256、关键版本和原归档内的证据入口。公开笔记不附鉴权信息、完整模型内部推理、私有运行路径或原始请求日志。

归档哈希用于标识复核对象，并不代替数学证明或独立外部认证。v04 离线修复小包中的“新增 POST 为 0”“候选未切换”等结论与该次前后封存指纹一致；v04c 的新增 6 次请求另有完整请求、SSE 终态和前后账本交叉核对。v04c 包的 1751 份清单文件均匹配。本文没有实时连接运行机器作网络或进程观测；追加本轮结果也没有覆盖原 v04 被截断的历史。

v05 完整过程包的 3296 份清单文件均与导出哈希及字节长度一致。12 次请求有完整终态证据；此前 1032 笔账本逐值保留，其中 79 笔未知不变。本次审阅只读取封存数据并执行自编的纯整数数学检查，没有执行包内运行器、调用研究模型或连接用户的本地现场。源码与工程候选的导出脱敏差异和原始指纹分别核对，不把导出字节误当原私有字节。

## v06 初始设计：明确提供 Sol 同伴，观察实际合作

用户已授权尝试加入 GPT-6.1 Sol。下一轮保留原 Adaption 负责人，初始提供一个 Sol/high 节点，由原负责人选择具体合作问题，经现有诉求、工作正文和版本授权后派发；第二名仍需明确缺口与方案。初始同伴记为 host_seeded。此次目标是观察有同伴时如何形成分工、质疑和整合，不把外部提供同伴算作自发创建。

静态源码审阅发现 v05 不能只改模型名称就获得同等边界：Adaption 的继承规则与 Gate 都要求相同 profile；现有费用守卫只计量 Adaption POST，监督器只扫描该渠道日志。Codex 原生 turn/start 内部可能包含多次模型交互，不能记成一次底层 POST 或套用 Adaption Credits。设计改为分渠道记录资源、共享持久停止；原生模型目录与实际确认仍在现场核对。

另一处会直接阻断联合工作：负责人无工具调用的最终响应会触发全组停止。如果派发后用普通最终回复表示“等待”，同伴会被清理。建议添加只读、有界 await_helper：宿主等待同一获准任务，不启动新模型、不给新额度，让结果返回仍运行的负责人工具循环。真正 final 仍结束阶段，避免无限自动唤醒。

保守实施边界为新增 Adaption 16 POST、最多 3.0408704 Credits 最坏责任；所有 Sol 共用 4 次原生 turn，每 turn 本地最多 300 秒、12 次工具，累计 48 次工具；首模型准入起阶段最多 45 分钟。数值为落实用户联研授权选定的实施限制，不冒充用户逐项原话。两种渠道分别计量，Codex 实扣与内部请求数没有被假定可见；不得自动增加付费渠道或越额。历史实际 38 次和 79 笔未知保留，旧 v05 的 42 上限与关闭状态不改，新阶段 Adaption 同根上限为 54。

局部次数到达上限只拒绝该资源后续准入，已准入执行与预留的负责人接收/整合在原总额内完成；第 4 次 Sol turn 不应刚准入就被取消。新未知、期限、无法继续的异常或负责人真正结束才触发全组停止。新同伴 turn 启动前预留至少两次负责人 POST 用于后续交付与整合，等待本身不做付费轮询。

以上是首次交接时的设计与静态审阅，当时尚未实现 v06 候选或启动模型；之后的实施与两次启动失败见下节。交给持有本机运行环境的 Codex 的接续包要求完成最小接入及针对性验证后执行一次有界试跑，不重复索取已经给出的 Sol 加入授权。后续报告须分别证明提议、运行、实际交付、采用或反驳，记录未合作、未送达与被截断结果。

外部研究者另整理了[固定进制有限深度的筛选补充](https://github.com/epsilon-lain/lain-s-note/blob/main/research/finite-base-observation-2026-11-11/FIXED_DEPTH.md)。它不是 v06 研究成果，数值范围也不超过已有笔记；不默认作为研究者初始解法。若实际检索或交付，需标记外部来源及后续采用。

## 观察十一：请求前的集成失败不能用来判断模型的合作选择

v06 首次启动在创建同伴时被旧的全局 profile 校验拒绝。隔离修复后，用户明确批准了一次合计不加额的人工接续。v06c 已创建一个宿主预置的 Sol/high 节点，但其任务数为 0，尚无 session。原负责人随后通过正常入口建立新上下文，仍在首个模型 POST 之前失败，原始错误为 `Mixed lead tool contract omits bounded helper await`。因此，“未联合研究”在这里来自执行入口阻塞，不能解释成模型不愿合作、没有分工能力或拒绝使用同伴。

根因是本地工具名 `ai_await_helper` 经真实请求转换后为 `yohaku_ai_await_helper`，入口却仍查旧名称。未激活的修复还处理了一个相邻缺口：现有本地参数校验器不支持该声明使用的 `number`，于是将等待秒数明确为 0–300 的整数；supervisor 在初始化和修改上下文前校验实际请求工具契约。共改动三个运行模块，修复候选 manifest SHA-256 为 `dec4a379e5886af6397b4d7b029c2906d98af04dd950827c1c68241b6e9c6a86`。

归档中的完整入口回归实际经过 Control、Rust 派发、run_task、Adaption.run、工具转换、参数校验和负责人工具循环；目录与推理 HTTP 使用替代实现。旧候选重现原错，修复候选两项检查通过，包括模拟同伴结果进入负责人的下一轮请求、缺少等待声明时在初始化前拒绝。外部审阅另核对七个报告源码指纹，并用提取的纯函数完成八项无网络检查，没有重跑包内控制器。

这比只在子模块或单独原生入口测试更接近真实失败点，但辅助者完成状态和返回值由夹具置入，等待为零秒，HTTP 为模拟 JSON 响应。它没有验证真实 Sol 执行、进行中的等待或供应商 SSE 流，更没有证明合作产物已被采用。旧的分段测试与本次完整入口覆盖之间的差异，是工程观察，不是新的模型能力结果。

## 观察十二：未花额度的关闭阶段，仍需要保留接续谱系

独立比较表明，前后 1044 条预留记录逐值相同，含原 79 笔未知；累计最坏责任值未变。两次失败合计新增 Adaption POST 0、Sol 原生 turn 0、费用预留 0。已有一个同伴占用 1/2 名额；节点创建与真实研究执行分别计数。原有 38 次 Adaption POST 保留。共享上限仍为新增 16 POST、3.0408704 Credits 最坏责任、4 次原生 turn 和 48 次原生工具，Adaption 同根最多 54 次。这里的 Credits 为责任约束，不是供应商账单审计。

两阶段均已关闭，持久停止记录保留。首次清理未完成也保留；后续封存快照支持相关运行与收尾进程退出、共享租约释放，不能将“阶段关闭”和“最后一个收尾进程退出”混为同一事件。首次模型准入与 deadline 均为空，45 分钟模型阶段时钟尚未启动。

额外静态审阅发现，现有 mixed gate 只登记原阶段和唯一 `continuation-01`，额度汇总也只覆盖两者；不可变 envelope 绑定了这一次接续的策略哈希。未激活的三模块修复没有修改这个限制。因此不能直接激活修复后换一个阶段 ID 启动，更不能删停止记录、覆盖旧策略或重建同伴以刷新名额。

下一次人工接续需先在同一共享额度下追加登记，显式纳入全部前驱阶段的请求、责任、原生 turn/工具、既有同伴和最早 deadline，并保留两次关闭记录、原负责人失败上下文及最后确认的模型响应。当前授权正文限定“一次人工接续”，已用于 v06c；本审阅不替用户扩写为无限重试。该迁移尚未实现或通过本包验证，检查点本身也没有派发权限。

本次 ZIP 的 3414 份清单内容均与导出哈希及字节长度一致。外部审阅只检查档案、源码和无网络纯函数，未连接用户现场或调用研究模型。本轮没有新增数学成果；已有外部数学反馈尚不能因放入包内就记作模型收到或采用。
````

## 输出与关联
```yaml
output: "research/pi-workflow-2026-11-12/README.md"
```
